% Generated draft: requires mathematical and source review.
\documentclass[11pt]{article}
\usepackage[margin=25mm]{geometry}
\usepackage{fontspec}
\setmainfont{Latin Modern Roman}
\usepackage{amsmath,amssymb,mathtools,unicode-math}
\setmathfont{Latin Modern Math}
\usepackage{enumitem,needspace,xurl,hyperref}
\setlength{\parindent}{0pt}
\setlength{\parskip}{5pt}
\setlength{\emergencystretch}{3em}
\newcommand{\N}{\mathbb{N}}
\newcommand{\Z}{\mathbb{Z}}
\newcommand{\Q}{\mathbb{Q}}
\newcommand{\R}{\mathbb{R}}
\newcommand{\C}{\mathbb{C}}
\newcommand{\F}{\mathbb{F}}
\newcommand{\A}{\mathbb{A}}
\newcommand{\PP}{\mathbb{P}}
\newcommand{\E}{\mathbb{E}}
\newcommand{\eps}{\varepsilon}
\newcommand{\dd}{\,\mathrm{d}}
\begin{document}
\section*{UnsolvedMath ID: 30005965}
\textbf{Source problem code:} catalogue entry
\subsection{Definitions and mathematical statement}
The supplied record does not determine a standalone mathematical conjecture. Its notation can be described only schematically: \(X\) is the space under discussion, \(G\) is the proposed symmetry group, \(\operatorname{Fuk}_{G}(X)\) denotes a sought-for \(G\)-equivariant Fukaya category, \(HH^{\bullet}\) denotes Hochschild cohomology, and \(QH^{\bullet}_{G}(X)\) denotes the intended equivariant quantum cohomology. The hypotheses on \(X\), \(G\), and their interaction, as well as coefficients and the relevant constructions, are absent from the supplied excerpt.

The source asks three linked questions:
\begin{enumerate}
\item Can one construct a suitable \(G\)-equivariant Fukaya category \(\operatorname{Fuk}_{G}(X)\)?
\item If such a category exists, does an identification
\[
HH^{\bullet}\!\left(\operatorname{Fuk}_{G}(X)\right)
\stackrel{?}{\cong} QH^{\bullet}_{G}(X)
\]
hold? The displayed isomorphism is a schematic rendering of the source's word ``equal''; the intended comparison and the structures it should preserve are unspecified.
\item What is the corresponding mirror structure?
\end{enumerate}
These are questions about the surrounding mathematical framework, with no recoverable quantifier over a specified class of spaces or group actions. They must not be interpreted as an assertion that the displayed identification holds universally.

The supplied statement-verification note says that the report uses these questions to motivate a talk, immediately rejects the naive first two formulations, and develops a more elaborate categorical framework. It explicitly says that the passage does not present a standalone unresolved conjecture. The supplied dated literature assessment further reports negative or more complicated answers in the \(U(1)\) case and describes the broader multipart question as unsettled. These are supplied source assessments, not independently verified conclusions.

Two apparent footnotes are appended to the excerpt. Footnote 2 says that an unidentified space is actually a spectrum constructed from Floer homotopy theory. Footnote 3 says that \(\check X\) is the dual torus to the torus in \(X\), and that a one-parameter subgroup \(\sigma\) corresponds to a character \(\check\sigma\). Their referents and role in the three questions are not supplied, so they cannot be imposed as precise hypotheses or used to specify the mirror structure.
\subsection{Short English statement}
The passage asks for an equivariant Fukaya category, a possible comparison of its Hochschild cohomology with equivariant quantum cohomology, and a mirror description. The supplied record marks the statement as unrecoverable and says the naive formulations are rejected in the report. Missing context prevents a precise standalone conjecture; this is a statement-only draft requiring human mathematical review.
\subsection{Sources}
\begin{itemize}
\item Supplied local UnsolvedMath registry record, permanent ID \(30005965\), title ``Equivariant Fukaya-Cohomology Discussion,'' source problem code \texttt{catalogue entry}. The project identifies \texttt{data/unsolvedmath/problems.json} and the supplementary \texttt{catalog/problems.json} as local registry files; their contents were not independently inspected for this draft.
\item Supplied report citation: \emph{Algebraic Geometry: Wall Crossing and Moduli Spaces, Varieties and Derived Categories} (2025), Oberwolfach Reports, DOI \texttt{10.4171/owr/2024/32}, pp.\ 1787--1854. Supplied source URL and literature-source link, labelled ``original OWR report'': \url{https://doi.org/10.4171/owr/2024/32}. The record gives proposed year 2024 and supplies no named proposer or talk author.
\item Literature assessment dated exactly \texttt{2026-08-22}: the supplied status is ``partially solved.'' The assessment reports negative or more complicated answers to the proposed equivariant Fukaya and Hochschild-equals-quantum-cohomology formulations for \(U(1)\), while describing the broad multipart question as unsettled. This is supplied dated curation, not independently verified literature assessment. The background explicitly requests independent verification of time-sensitive status and cited claims.
\item Supplied statement status: ``unrecoverable.'' The supplied statement-verification note says that the report rejects the naive first two questions and develops a more elaborate categorical framework, without presenting a standalone unresolved conjecture. The report was not independently visited or checked for this draft.
\end{itemize}
\end{document}
